\documentclass[11pt,a4paper]{article}
\usepackage[margin=20mm]{geometry}
\usepackage[T1]{fontenc}
\usepackage{lmodern}
\usepackage{amsmath,amssymb}
\usepackage{array,booktabs,longtable}
\usepackage[hidelinks]{hyperref}
\usepackage{url}
\urlstyle{tt}
\newcommand{\code}[1]{\ifmmode\text{\nolinkurl{#1}}\else\nolinkurl{#1}\fi}
\newcolumntype{L}[1]{>{\raggedright\arraybackslash}p{#1}}
\newenvironment{vartable}{%
  \small
  \begin{longtable}{@{}L{0.22\textwidth}L{0.49\textwidth}L{0.23\textwidth}@{}}
  \toprule
  Variable & Meaning & Units, shape, or source \\
  \midrule
  \endhead
  \bottomrule
  \endfoot
}{\end{longtable}}
\setlength{\emergencystretch}{2em}
\renewcommand{\arraystretch}{1.18}

\title{PyKlyC: Post-processing Variables and Display Parameters}
\author{}
\date{}

\begin{document}
\maketitle

This reference covers the implemented diagnostics following the nonlinear
iteration and the first two plots: the Applegate diagram and current
harmonic amplitudes. It follows \code{KLJCmainsv22.m}, particularly lines
640--773. The final section documents all six entries of MATLAB's
\code{GIFpara} vector, including their later animation uses.

Unless stated otherwise, Python indices are zero-based. Let
\(N_z^{\mathrm{grid}}=\operatorname{len}(\code{grid.z})\) and
\(H=\code{Ntimes}\,N_z^{\mathrm{grid}}\).
These are explanatory shape abbreviations, not new program variables.
The quantity \code{Nz} is a mesh-resolution setting; it is not generally
equal to \(N_z^{\mathrm{grid}}\).

\section{Reference quantities and coordinates}

\begin{vartable}
\code{U0} & Reference beam voltage. & kV \\
\code{I0} & Magnitude of the reference average current, per beam. & A \\
\code{Iz0} & Signed reference average current. & A \\
\code{Nbeam} & Number of beams. & Dimensionless \\
\code{s} & Reference velocity divided by the speed of light: \(v_0/c\). & Dimensionless \\
\code{betae} & Reference electronic propagation constant, \(\omega/v_0\), expressed for millimetre coordinates. & \(\mathrm{mm}^{-1}\) \\
\code{zi} & Physical longitudinal grid. & \((N_z^{\mathrm{grid}},)\), mm \\
\code{z} & Normalized longitudinal grid: \(z=\beta_e z_i\). & \((N_z^{\mathrm{grid}},)\), dimensionless \\
\code{dz} & Physical longitudinal grid spacing. & mm \\
\code{grid.dt} & Normalized step \(\beta_e\,dz\); also the phase increment used in time tracking. & Dimensionless \\
\code{esign} & Charge-sign convention used in the force and energy equations. & \(+1\) for electrons \\
\code{eta0} & Vacuum wave impedance. & \(\Omega\) \\
\(P_{\mathrm{beam}}\) & Total reference beam power: \(N_{\mathrm{beam}}(1000U_0)I_0\). & W; explanatory notation \\
\end{vartable}

\section{Beam and particle quantities read by the diagnostics}

\begin{vartable}
\code{Nt} & Number of emitted samples over one RF period. & Integer \\
\code{Ns} & Number of retained RF harmonics. & Integer \\
\code{Ntimes} & Multiplier allocating trajectory-history capacity. It is not a number of RF periods. & Integer \\
\code{Y} & Active particle-column indices. & Integer array \\
\code{Yflag} & Number of particles counted as exiting downstream. & Integer counter \\
\code{tn1} & Arrival phase of each particle at the entrance. & \((Nt,)\), radians \\
\code{Tn} & Phase/time history \(\omega t\) along tracked trajectories. & \((H,Nt)\) \\
\code{Zn} & Normalized position history along tracked trajectories. & \((H,Nt)\) \\
\code{Vn} & Signed velocity history, \(v/v_0\). & \((H,Nt)\) \\
\code{Tnz} & Arrival phase at longitudinal grid-plane crossings. & \((N_z^{\mathrm{grid}},Nt)\) \\
\code{IzDC} & DC current divided by \code{Iz0}. & \((N_z^{\mathrm{grid}},)\) \\
\code{Iz} & Complex current-harmonic amplitudes divided by \code{Iz0}; column \(k\) in Python represents harmonic \(k+1\). & \((N_z^{\mathrm{grid}},Ns)\) \\
\code{effe} & Beam kinetic-power reduction relative to reference input beam power, evaluated along the tube. & \((N_z^{\mathrm{grid}},)\), dimensionless \\
\end{vartable}

\code{Tn} and \code{Zn} describe the trajectory history, including any
continuation by time tracking. \code{Tnz} describes grid-plane crossings.
Unwritten entries of \code{Zn} and \code{Vn} are positive infinity and
must be excluded when plotting. The arrival phases are stored unwrapped;
phase wrapping is a display operation.

\section{Cavity and field quantities}

\begin{vartable}
\code{Nc} & Number of cavities or interaction structures. & Integer \\
\code{type} & Input: 0; intermediate: 1; output: \(-1\); traveling-wave: 2. & \((Nc,)\) \\
\code{nh} & Harmonic associated with each cavity; zero denotes a static gap. & \((Nc,)\) \\
\code{zshift} & Cavity-center positions. & \((Nc,)\), mm \\
\code{k0} & Resonant free-space wavenumber of each cavity. & \((Nc,)\), \(\mathrm{mm}^{-1}\) \\
\code{Qin} & Intrinsic cavity quality factor. & \((Nc,)\), dimensionless \\
\code{gaman} & Output-load reflection coefficient. & \((Nc,)\), complex \\
\code{Pin} & Nominal incident RF power. & \((Nc,)\), W \\
\code{Vin} & Normalized incident-wave amplitude before ramp scaling. & \((Nc,)\) \\
\code{drivefactor} & Incident-voltage multiplier; its square multiplies incident power. & Scalar \\
\code{Vnc} & Normalized complex cavity-mode amplitude. & \((Nc,)\) \\
\code{Vnp} & Normalized port-wave amplitude: reflected wave at an input, outgoing wave at an output. & \((Nc,)\), complex \\
\code{ez} & Cavity profiles sampled on \code{zi}. RF profiles have units \(\mathrm{mm}^{-3/2}\); static profiles use kV/mm. & \((N_z^{\mathrm{grid}},Nc)\) \\
\code{Esz0} & Separate normalized DC space-charge field in Python. & \((N_z^{\mathrm{grid}},)\) \\
\code{Esz} & Separate normalized RF space-charge harmonics. & \((N_z^{\mathrm{grid}},Ns)\), complex \\
\end{vartable}

The normalized profile \code{ez} is distinct from the cavity contribution
\code{Emz} used by tracking. The total RF tracking field is \code{Esmz};
the corresponding DC names are \code{Emz0} and \code{Esmz0}.
MATLAB reuses \code{Esz0} for its combined DC field. The implemented
space-charge diagnostic reads the separate space-charge field; these agree
in the currently supported solver because \code{Emz0} is zero.

\section{Implemented diagnostic outputs}

\subsection{Trajectory and modulation diagnostics}
All functions in this table belong to \code{src/diagnostics/beam.py}.

\begin{vartable}
\code{n} & Largest last-written history row among active particles. Retains MATLAB's one-based row numbering. & \code{check_boundary_reaching()} \\
\code{status} & MATLAB bookkeeping classification: \code{Storage limit}, \code{Reflected electrons}, or \code{Normal}. & \code{check_boundary_reaching()} \\
\code{minimum_v_over_c} & Minimum recorded \code{Vn} multiplied by \code{s}. & \code{check_negative_velocity()} \\
\code{negative_velocity} & Whether the minimum recorded velocity is negative. & Boolean; \code{check_negative_velocity()} \\
\code{Izmax} & Maximum of \code{abs(Iz[:, 0])}. & \code{find_peak_modulation()} \\
\code{ntemp} & Grid index of strongest fundamental current modulation. & Zero-based; \code{find_peak_modulation()} \\
\code{ncenter} & Particle-column index nearest the fundamental current-peak phase at that plane. & Zero-based; \code{find_peak_modulation()} \\
\end{vartable}

The \code{Reflected electrons} status means missing downstream exits under
MATLAB's bookkeeping; actual recorded negative velocity is checked separately.
MATLAB reuses temporary names: \code{ntemp} inside the boundary check is a
last-written history row, whereas the modulation diagnostic returns a grid
index with that name.

The fundamental current varies as
\[
\operatorname{Im}\!\left(Iz_1 e^{i\theta}\right)
=|Iz_1|\sin\!\left(\theta+\arg Iz_1\right).
\]
Its peak phase is \(\theta_{\mathrm{peak}}=\pi/2-\arg Iz_1\), modulo
\(2\pi\). The selected particle minimizes the distance to this phase on
the unit circle. Higher harmonics can shift the actual current peak;
\code{ncenter} is a plotting reference, not an exact bunch centroid.

\subsection{Cavity-power diagnostics}
These outputs belong to \code{calculate_cavity_power()} in
\code{src/diagnostics/cavity.py}.

\begin{vartable}
\code{Vng} & Complex gap voltage for each cavity. & \((Nc,)\), kV \\
\code{effloss} & Each cavity's wall-loss power divided by total reference beam power. & \((Nc,)\), dimensionless \\
\code{lossacc} & Cumulative wall-loss fraction along the tube; includes half a cavity's loss exactly at its center. & \((N_z^{\mathrm{grid}},)\) \\
\code{RFinput} & Total net RF input power after reflection. & W \\
\code{refcavity} & Complex reflection coefficient of the last input cavity. & Scalar, complex \\
\code{effcavity} & Load power from the last output cavity divided by total reference beam power. & Scalar, dimensionless \\
\end{vartable}

\subsection{Space-charge exchange and final power balance}

\begin{vartable}
\code{effSCDC} & Accumulated DC space-charge energy-exchange contribution relative to the entrance. & \code{calculate_space_charge_power()} \\
\code{effSCRF} & Accumulated RF space-charge energy-exchange contribution relative to the entrance. & \code{calculate_space_charge_power()} \\
\code{effSC} & Sum \code{effSCDC + effSCRF}. & \code{calculate_space_charge_power()} \\
\code{effend} & Beam efficiency sampled at the observation position. & \code{calculate_power_balance()} \\
\code{effSCend} & \code{effSC} sampled at the same position. & \code{calculate_power_balance()} \\
\code{effendmodified} & RF efficiency inferred from the final energy balance. & \code{calculate_power_balance()} \\
\end{vartable}

The first function is in \code{src/diagnostics/space_charge.py}; the second
is in \code{src/diagnostics/power_balance.py}. All quantities beginning
with \code{eff} are fractions, not percentages. The accumulated space-charge
terms can have either sign; they are not necessarily dissipative losses.
With \(\Delta z_{\mathrm{norm}}=\beta_e\,dz\), the left-endpoint updates are
\begin{align}
\mathrm{effSCDC}_j &= \mathrm{effSCDC}_{j-1}
 + e_{\mathrm{sign}}\Delta z_{\mathrm{norm}} IzDC_{j-1}Esz0_{j-1},\\
\mathrm{effSCRF}_j &= \mathrm{effSCRF}_{j-1}
 + \frac{e_{\mathrm{sign}}\Delta z_{\mathrm{norm}}}{2}
 \sum_{m=1}^{N_s}\operatorname{Re}\!\left(Iz_{j-1,m}^{*}Esz_{j-1,m}\right).
\end{align}
The final balance is
\[
\boxed{\mathrm{effendmodified}
 =\mathrm{effend}-\mathrm{effSCend}
 -\sum_k\mathrm{effloss}_k
 +\frac{\mathrm{RFinput}}{P_{\mathrm{beam}}}.}
\]

\section{Quantities for the first two plots}

These plotting functions are implemented in \code{src/plotting/beam.py} as
\code{plot_applegate()} and \code{plot_current_harmonics()}.
The symbols \(\psi\), \(\psi_c\), and \(\psi_{\mathrm{plot}}\) below
explain MATLAB's expression rather than introduce existing state variables.

\begin{vartable}
\code{zpeak} & \code{grid.zi[ntemp]}: physical position of strongest fundamental modulation. & mm \\
\code{phase_reference} & \code{beam.tn1[ncenter]}: selected electron's entrance phase. & radians \\
\code{Ntp} & MATLAB display parameter used to set trajectory-sampling density; also controls animation time sampling. & Positive integer \\
\code{kp} & Column index of a particle being plotted. & Integer \\
\code{Zn[:, kp] / betae} & Physical particle position. & Applegate horizontal axis, mm \\
\(\psi=\code{Tn}-\code{Zn}\) & Phase relative to motion at the reference velocity: \(\omega t-\omega x/v_0\). & Unwrapped phase \\
\(\psi_c=\code{tn1[ncenter]}-\code{z[0]}\) & Center of the displayed phase window. & radians \\
\(\psi_{\mathrm{plot}}\) & Phase wrapped into \([\psi_c-\pi,\psi_c+\pi)\). & Applegate vertical axis \\
\code{ez_background} & Magnitude of the sum of RF cavity profiles, suppressing static-gap columns. & Second Applegate axis, \(\mathrm{mm}^{-3/2}\) \\
\code{abs(Iz[:, k])} & Normalized amplitude of harmonic number \(k+1\) in Python. & Current-harmonic vertical axis \\
\code{zmax} & Right plotting limit; proposed usage is \code{max(zi)} minus \code{efficiency_end_offset_mm}. & mm \\
\end{vartable}

The Applegate phase wrapping is
\[
\psi_{\mathrm{plot}}=\psi-2\pi
\left\lfloor\frac{\psi-\psi_c+\pi}{2\pi}\right\rfloor.
\]
Its field background is
\[
\mathrm{ez\_background}(z_i)
=\left|\sum_k ez(z_i,k)\operatorname{sign}(nh_k)\right|.
\]
This is a normalized profile, not a field scaled by \code{Vnc} and not the
instantaneous force field experienced by the particles. Python should form
a separate background array without modifying \code{cavity_fields.ez}.

\section{MATLAB GIFpara and the Python mapping}

The exact MATLAB variable name is \code{GIFpara}, not \code{GIFparams}.
It is a six-entry vector mixing display parameters with quantities also
used in numerical calculations. MATLAB entry \code{GIFpara(j)} corresponds
to Python entry \code{GIFpara[j-1]}.

\begin{longtable}{@{}L{0.16\textwidth}L{0.38\textwidth}L{0.40\textwidth}@{}}
\toprule
Entry & Meaning and MATLAB use & Python counterpart / implementation status \\
\midrule
\endhead
\bottomrule
\endfoot
\code{GIFpara(1)}\newline Python \code{[0]}
& Sets \code{Ntp}. Controls displayed-particle subsampling in trajectory plots.
  For animations, sets the normalized time step \(dtp=2\pi/Ntp\), i.e.
  \(Ntp\) time intervals per RF period. Dimensionless positive integer.
& Stored in \code{DisplaySettings.GIFpara[0]}. Plotting argument:
  \code{Ntp}. Full animation sampling is not yet implemented. It does not
  set the solver's \code{Nt} or \code{grid.dt}. \\
\code{GIFpara(2)}\newline Python \code{[1]}
& Sets \code{ymin}, the lower velocity-axis limit. Applies to either
  \(v/c\) or \(\gamma v/c\), according to entry 6. Dimensionless.
& Stored in \code{DisplaySettings.GIFpara[1]}. A display limit; no effect
  on forces, velocities, or convergence. Not used for the first two plots. \\
\code{GIFpara(3)}\newline Python \code{[2]}
& Sets \code{ymax}, the upper velocity-axis limit, using the same convention
  as entry 2. Dimensionless.
& Stored in \code{DisplaySettings.GIFpara[2]}. A display limit; no effect
  on the computed solution. Not used for the first two plots. \\
\code{GIFpara(4)}\newline Python \code{[3]}
& Upstream cut distance in mm: \(z_{\min}=\min(zi)+GIFpara(4)\).
  Sets the animation's lower position limit. MATLAB also discards a
  backward-moving particle when the lower crossed grid plane is below this
  cutoff (line 503).
& Tracking uses the separate numerical setting
  \code{NumericalSettings.upstream_discard_distance_mm}.
  \code{DisplaySettings.GIFpara[3]} is the display cut. When exporting
  tracking inputs to MATLAB, use the numerical discard distance in entry 4.
  This value is a distance, not an absolute position. \\
\code{GIFpara(5)}\newline Python \code{[4]}
& Downstream cut distance in mm:
  \(z_{\max}=\max(zi)-GIFpara(5)\).
  Also sets the efficiency observation position for convergence and power
  balance. Later selects the nearest grid plane for spent-beam analysis.
& Numerical calculations use
  \code{NumericalSettings.efficiency_end_offset_mm}.
  \code{DisplaySettings.GIFpara[4]} is a separate display cut. The proposed
  first two plots can receive an explicit \code{zmax}; use the numerical
  observation position when matching MATLAB. Export that numerical offset
  as entry 5 for MATLAB comparisons. \\
\code{GIFpara(6)}\newline Python \code{[5]}
& Sets \code{velocitytype}. Zero plots \(v/c\); any nonzero value plots
  \(\gamma v/c\), where \(\gamma=[1-(v/c)^2]^{-1/2}\).
  The latter is longitudinal momentum divided by \(mc\) in this 1D model.
& Stored in \code{DisplaySettings.GIFpara[5]}. Selects a future velocity
  or phase-space plot's vertical quantity; it does not change tracking.
  Not used for the first two plots. \\
\end{longtable}

\subsection{Details needed for a faithful translation}

\begin{itemize}
\item The Applegate diagram chooses a stride
  \(\max(\operatorname{round}(Nt/Ntp),1)\), takes the corresponding emitted
  sample indices, and intersects them with \code{Y}. Thus \code{Ntp} is not
  a guarantee of exactly that many plotted trajectories. The later MATLAB
  velocity plot uses \(\operatorname{round}(Nt/Ntp)\) without the lower bound.
\item MATLAB's positive half-integer rounding differs from NumPy's default
  ties-to-even rounding. Exact subsampling parity should account for this.
\item The first two MATLAB figures use \([\min(zi),z_{\max}]\), not
  \([z_{\min},z_{\max}]\). Entry 4 therefore does not crop their left edge,
  although it can already have affected reflected-particle tracking.
\item Entries 2, 3, and 6 concern velocity displays, not the Applegate
  phase axis or the current-harmonic amplitudes.
\item The final efficiency uses interpolation at \(z_{\max}\). Spent-beam
  analysis instead rounds to a grid-plane index; these locations need not
  coincide exactly when the offset is not an integer number of mesh steps.
\item \code{GIFpara} specifies parameters. Figure and movie generation are
  enabled by separate controls such as \code{figuredisp} and \code{moviedisp}.
\item \code{DisplaySettings} already stores the legacy vector, but this does
  not mean all of its plotting and animation consumers exist in Python yet.
\end{itemize}

\section{Source locations}
\begin{itemize}
\item \code{KlyC/KLJCmainsv22.m}: line 503 (upstream discard), line 532
  (efficiency observation), lines 640--724 (diagnostics), lines 729--773
  (first two plots), lines 799--820 (velocity plot), line 840 (spent-beam
  plane), and line 970 onward (animation sampling).
\item \code{pyKlyC/src/model/settings.py}: \code{NumericalSettings} and
  \code{DisplaySettings} definitions.
\item \code{pyKlyC/src/diagnostics/}: implemented trajectory, modulation,
  cavity-power, space-charge-power, and final power-balance functions.
\end{itemize}

\end{document}
